\subsection{Empirical evidence argues against symmetric structural equivalence}

The directional witness is not only a conservative formal choice. Wang reports a controlled asymmetry in structural transfer between natural-language and biological foundation models \cite{wang2026asymmetry}. Language models trained or fine-tuned on structural language tasks retain meaningful competence when moved to biological-sequence structure, while the reverse transfer remains weak across matched-token controls, scaling and multiple biological model families. The result is domain-specific and does not establish a universal law of transfer, but it is a direct counterexample to the assumption that shared structural regularity implies symmetric representational transfer.

For Paper 3 this has two consequences. First, a structural relation should default to a directed edge $S_A\to S_B$ with a witness rather than an undirected equivalence class. Second, transfer evaluation should include direction reversal as an explicit perturbation. If $A\to B$ succeeds but $B\to A$ fails, the library should preserve that asymmetry rather than canonicalize the two tasks into one globally interchangeable class. Any later transitive composition likewise requires a new scope/invariant check.

This also strengthens the Q2/Q3 benchmark logic. The object being tested is not whether two domains ``look analogous'' in a human narrative sense. It is whether a particular source-side operator or inference depends only on relations and invariants that are actually available in the target under the registered boundary conditions. Directional negative-transfer history is therefore part of the reusable state.

\subsection{Executable transport obligations and fail-closed abstention}

The strengthened implementation makes the witness more specific than a structural-similarity score. A directional witness is interpreted as a finite set of typed transport obligations
\[
\Omega_{A\to B}^{(q)}=\{o_1,\ldots,o_m\},
\]
scoped to quantity of interest $q$. An obligation records a role, relation, invariant, source precondition, target boundary, QoI requirement or forbidden loss together with its requirement class, evidence identities and a three-valued verification status
\[
\operatorname{status}(o_i)\in\{\mathrm{SATISFIED},\mathrm{VIOLATED},\mathrm{UNKNOWN}\}.
\]
The distinction between \texttt{VIOLATED} and \texttt{UNKNOWN} is load-bearing. Failure to establish a required precondition is not itself evidence that the precondition is false.

The transfer decision is deliberately non-compensatory:
\[
\operatorname{Transfer}(A\to B;q)=
\begin{cases}
\mathrm{LICENSED}, & \text{all load-bearing obligations are satisfied},\\
\mathrm{REJECTED}, & \text{at least one load-bearing obligation is demonstrably violated},\\
\mathrm{CANNOT\_CHECK}, & \text{otherwise, if at least one load-bearing obligation is unresolved}.
\end{cases}
\]
Thus several preserved relations cannot compensate for a violated target boundary, wrong transfer direction or forbidden loss. Conversely, a property declared irrelevant to the registered QoI may be listed as a permissible loss rather than forcing global structural equivalence. Each load-bearing obligation is content-bound to evidence rather than inheriting authority merely because the witness object contains some evidence somewhere. The implementation emits an obligation-level trace and a deterministic content hash so that a confirmatory benchmark can freeze the exact witness contract before outcome access.

These semantics are an implementation and evaluation contract, not an empirical superiority result. The confirmatory question remains whether machine-produced obligations improve valid distant transfer and reduce hostile semantic near-misses relative to strong semantic, relational and mechanism-alignment controls under matched resources.
